\documentclass[a4paper,10pt]{article}
\usepackage[T1]{fontenc}
\usepackage[utf8]{inputenc}
\usepackage{lmodern,textcomp}
\usepackage[margin=2.2cm]{geometry}
\usepackage{amsmath,amssymb,siunitx,booktabs,graphicx,tabularx,array,enumitem,xcolor,fancyhdr}
\usepackage[colorlinks=true,linkcolor=blue!50!black,urlcolor=blue!50!black]{hyperref}
\definecolor{draft}{RGB}{135,70,0}
\setlength{\parindent}{0pt}
\setlength{\parskip}{4pt}
\raggedright
\setlength{\headheight}{14pt}
\renewcommand{\arraystretch}{1.2}
\renewcommand{\tabularxcolumn}[1]{>{\raggedright\arraybackslash}p{#1}}
\setlist{nosep,leftmargin=*}
\newcommand{\todo}[1]{\par\smallskip{\color{draft}\textbf{To complete:} #1}\par\smallskip}
\newcommand{\blank}{\textcolor{draft}{\textit{[fill in]}}}
\newcommand{\file}[1]{\nolinkurl{#1}}
\newcommand{\fig}[3][0.92]{\begin{center}\includegraphics[width=#1\linewidth]{#2}\\{\small #3}\end{center}}
\pagestyle{fancy}\fancyhf{}
\fancyhead[L]{34870 Electroacoustics}
\fancyhead[R]{Lab B --- Microphone scattering}
\fancyfoot[L]{Group 10 \quad / \quad measured 22 September 2026}
\fancyfoot[R]{\thepage}
\graphicspath{{../figures/}{Lab B/figures/}}
\title{34870 Electroacoustics\\Lab B --- Microphone scattering}
\author{Group 10: Louis Andrianne, Sophie Kimura, Mads Rudolph}
\date{Laboratory date: 22 September 2026}
\begin{document}
\begin{center}
{\Large\bfseries Lab B --- Microphone scattering}\\[5pt]
34870 Electroacoustics\\[4pt]
Group 10: Louis Andrianne, Sophie Kimura, Mads Rudolph\\[4pt]
Laboratory date: Tuesday 22 September 2026, 10:00--12:00 (building 354, rooms 028/025)
\end{center}\medskip
\thispagestyle{fancy}
\textbf{Working draft.} Measurements taken and processed on 22 September 2026. The
coloured prompts mark what the group still has to decide or add (setup sketch,
handbook comparison, uncertainty budget). Raw data is unchanged; the analysed files
are listed in the processing record at the end.

\section{Purpose and theoretical basis}
The experiment investigates how a microphone-shaped rigid body changes the pressure
at its front face. A wooden, enlarged microphone model makes this scattering measurable
using a small microphone. The angle-dependent response is compared with a boundary
element method (BEM) calculation and with microphone handbook free-field corrections.

For a body radius $a$, the dimensionless frequency is
\begin{equation}
ka=\frac{2\pi f a}{c},\qquad a=\frac{D_{\mathrm{mock}}}{2}.
\end{equation}
Here $c$ is the sound speed. Similar geometries are compared at equal $ka$.
The measured voltage ratio and the normalized scattering response are
\begin{align}
H_{\theta}(f)&=\frac{V_{\mathrm{mic},\theta}(f)}{V_{\mathrm{ref},\theta}(f)},&
H_0^{\mathrm{ref}}(f)&=\frac{V_{\mathrm{mic,no\ mock}}(f)}{V_{\mathrm{ref,no\ mock}}(f)},\\
H_{\mathrm{norm},\theta}(f)&=\frac{H_\theta(f)}{H_0^{\mathrm{ref}}(f)},&
\Delta L_\theta(f)&=20\log_{10}|H_{\mathrm{norm},\theta}(f)|.
\end{align}
The subscript $0$ in $H_0^{\mathrm{ref}}$ denotes the reference run, not zero-degree
incidence. Normalization removes the common measurement-chain response only when its
settings and the loudspeaker and UMIK positions remain unchanged.

For $ka\ll1$ the body is small compared with the wavelength, the wave flows around it
and $\Delta L_\theta\to0$ at every angle. As $ka$ approaches 1 the front face starts to
reflect; head-on the reflected wave doubles the pressure (\SI{+6}{dB} for an infinite
rigid wall) and the wave diffracted at the rim of the face reaches the centre from the
whole edge with the same delay, so the centre-of-face pressure exceeds \SI{+6}{dB} and
reaches about \SI{+10}{dB} near $ka\approx3$. At grazing incidence ($90^\circ$) the face
does not obstruct the wave and the increase stays within about \SI{2}{dB}. The
loudspeaker radiates a spherical wave; over the body it approximates a plane wave when
the source distance $r$ is large compared with the body length and when the level
follows $1/r$ (Part~0). With $r=\SI{2.8}{m}$ and a body of \SI{0.855}{m}, the incident
phase along the body is not exactly planar; the effect is small at the low $ka$ where
the comparison is made.

\section{Equipment and geometry}
\par\smallskip
\begin{tabularx}{\linewidth}{@{}lX@{}}\toprule
Item & Recorded identity or setting \\\midrule
PC / internal sound card & Lab desktop, internal card (\texttt{Line In} reference, \file{Speakers} output); course routine \file{meas_mag_spec2_SoundCard_LabB} called from \file{labB.m} \\
UMIK MiniDSP / USB microphone & S/N 708-0332, correction for $90^\circ$ incidence on the UMIK (it points up) \\
Loudspeaker / amplifier & Chamber loudspeaker; amplifier setting not noted, \textbf{unchanged} for the whole session \\
Wooden microphone mock-up & $D_{\mathrm{mock}}=\SI{250}{mm}$, length \SI{855}{mm} --- \textbf{nominal values from the brief, not measured with the tape} \\
Mock-up back end & Round \\
Loudspeaker--UMIK distance & $r=\SI{2.8}{m}$ from the loudspeaker front, confirmed by the $1/r$ check; driver alignment and height not noted \\
UMIK tip--front-face gap & Approximately \SIrange{2}{3}{cm}, not flush (confirmed by the data, Sec.~\ref{sec:tip}); mock-up rotated around the fixed UMIK \\
Angle reference / temperature & Paper protractor under the mock-up; temperature not measured, $c=\SI{344}{m/s}$ assumed \\\bottomrule
\end{tabularx}\par\smallskip
Geometry as recorded in \file{Lab B/data/labB_log.md}. The lab PC clock runs ahead of
real time (files stamped 12:31--13:11 for a slot that ended at 12:00); the order of
the runs is what matters and is preserved.
\todo{Add the setup sketch (brief Fig.~1) with the fixed positions, the \SIrange{2}{3}{cm} gap and the mock-up angle; add the mock-up dimensions if somebody measures them in a later slot.}

\newpage
\section{Part 0: Anechoic chamber and free-field check}
The brief specifies a chamber cutoff frequency of \SI{125}{\hertz}. The expected
pressure magnitude of an approximately spherical wave varies as $1/r$, giving
\begin{equation}
\Delta L_{21}=20\log_{10}\!\left(\frac{|p(r_2)|}{|p(r_1)|}\right)
\simeq20\log_{10}\!\left(\frac{r_1}{r_2}\right).
\end{equation}
The distance check was carried out \emph{after} the scattering series, without the
mock-up: the UMIK was moved from its series position (\SI{2.8}{m}) to \SI{1.4}{m} and
\SI{0.93}{m}. Doing it last guaranteed that the microphone did not move between the
reference and the seven angles.
\par\smallskip
\begin{tabularx}{\linewidth}{@{}lXXX@{}}\toprule
Run & Distance / m & Data file & Observed difference / dB (mean \SIrange{125}{200}{Hz}) \\\midrule
Reference distance & 0.93 & \file{no_mockup_0.93.mat} $\to$ \file{labB_dist93cm.mat} & Reference \\
Second distance & 1.4 & \file{no_mockup_1.4m.mat} $\to$ \file{labB_dist140cm.mat} & $-3.66$ (predicted $-3.55$) \\
Third & 2.8 & \file{no_mockup_1.8m.mat} $\to$ \file{labB_dist280cm.mat} & $-9.61$ (predicted $-9.57$) \\\bottomrule
\end{tabularx}\par\smallskip
The third file is named \file{no_mockup_1.8m} on the lab PC; the tape reading was
\SI{2.8}{m}, and the observed $-9.6$~dB matches \SI{2.8}{m} ($-9.57$~dB), not
\SI{1.8}{m} ($-5.7$~dB).
\fig{labB_free_field_check.png}{Part 0: level relative to the \SI{0.93}{m} run, with the $1/r$ predictions (dashed).}
Between the chamber cutoff and about \SI{400}{Hz} both curves follow $1/r$ to within
\SI{0.1}{dB}: the chamber is free-field there and the loudspeaker radiates like a point
source. Above \SI{400}{Hz} the \SI{1.4}{m} curve ripples by $\pm3$--\SI{4}{dB}
(interference between the direct sound and reflections from the stand, the floor net or
the loudspeaker cabinet) and above \SI{4}{kHz} both curves wander by several dB (source
directivity; a \SI{1}{cm} placement error is already $0.1\lambda$). Below \SI{100}{Hz}
the \SI{2.8}{m} curve drops to \SI{-16}{dB}: below its cutoff the room reflects. The
useful free-field band is therefore \SIrange{125}{2000}{Hz} with a margin of about
\SI{1}{dB} at the higher frequencies.

\section{Part 1: Scattering measurement procedure}
\begin{enumerate}
\item The UMIK was connected before MATLAB was started; the course script found the
internal card (\texttt{Line In}, \texttt{Speakers}) and the UMIK. A trial run
(\file{test.mat}) was used to set the level.
\item The loudspeaker and the upward-pointing UMIK were fixed at \SI{2.8}{m}; the
mock-up was placed with its face centre \SIrange{2}{3}{cm} from the UMIK tip and
rotated on a protractor for $0^\circ$, $30^\circ$, $60^\circ$, $90^\circ$, then
$45^\circ$, $15^\circ$ and $75^\circ$.
\item The mock-up was removed and the reference measured at the same UMIK position.
No reference was taken \emph{before} the series, so no drift check is available.
\item Part 0 (Sec.~3) was measured last. Data was copied to a group folder and
verified at home.
\end{enumerate}
\todo{Setup sketch: sound-card output to line-in reference and amplifier; amplifier to loudspeaker; UMIK to PC through USB.}

\par\smallskip
\begin{tabularx}{\linewidth}{@{}lXlX@{}}\toprule
Parameter & Value used & Parameter & Value used \\\midrule
$f_1$ & \SI{50}{\hertz} & $f_2$ & \SI{10000}{\hertz} \\
Tones/octave & 24 (184 tones) & Resolution & \SI{1}{\hertz} \\
Averages $N_{\mathrm{av}}$ & 32 & UMIK incidence & $90^\circ$ \\\bottomrule
\end{tabularx}\par\smallskip
All as prescribed. The UMIK incidence selects its correction and is \emph{not} the
mock-up angle. Departure from the planned procedure: the course script \file{labB.m}
was used instead of the group's wrapper, so each saved file holds only
\file{h = specn(:,2)./specn(:,1)}; the frequency vector was reconstructed from the
multitone generator with the same parameters (\file{import_group_files.m}), and the
full spectra needed for a noise-floor plot were not saved.

\newpage
\subsection{Measurement record and quality checks}
\par\smallskip
\begin{tabularx}{\linewidth}{@{}lXlX@{}}\toprule
Condition & File (lab PC $\to$ analysed) & Angle & Quality / notes \\\midrule
Trial & \file{test.mat} & -- & level check, not used \\
Mock-up & \file{mockup_0.mat} $\to$ \file{labB_ang000.mat} & $0^\circ$ & clean below \SI{2}{kHz}; notch at \SI{3.9}{kHz} (tip gap) \\
Mock-up & \file{mockup_15.mat} $\to$ \file{labB_ang015.mat} & $15^\circ$ & notch at \SI{2.5}{kHz} \\
Mock-up & \file{mockup_30.mat} $\to$ \file{labB_ang030.mat} & $30^\circ$ & notch at \SI{3.4}{kHz} \\
Mock-up & \file{mockup_45.mat} $\to$ \file{labB_ang045.mat} & $45^\circ$ & notch at \SI{3.9}{kHz} \\
Mock-up & \file{mockup_60.mat} $\to$ \file{labB_ang060.mat} & $60^\circ$ & notch at \SI{6.0}{kHz} \\
Mock-up & \file{mockup_75.mat} $\to$ \file{labB_ang075.mat} & $75^\circ$ & no deep notch \\
Mock-up & \file{mockup_90.mat} $\to$ \file{labB_ang090.mat} & $90^\circ$ & no deep notch \\
No mock-up & \file{no_mockup_1.8m.mat} $\to$ \file{labB_nomockup.mat} & -- & reference, \SI{2.8}{m}, taken after the series \\
End reference & not taken & -- & no drift check possible \\\bottomrule
\end{tabularx}\par\smallskip
No repeats were needed. Clipping cannot be excluded from the saved data (no full
spectra), but the ratios are smooth and consistent from angle to angle, and the trial
run was used to set a moderate level.

\subsection{Normalized angle responses}
\fig{labB_measured_vs_bem.png}{Measured $\Delta L_\theta$ (solid) and the BEM centre-of-face prediction (dashed) for the same angles.}
\par\smallskip
\begin{tabularx}{\linewidth}{@{}lXXXX@{}}\toprule
Incidence & Level \SIrange{125}{200}{Hz} / dB & Maximum below \SI{2}{kHz} & At \SI{1345}{Hz} / dB & BEM face centre \\\midrule
$0^\circ$ & $+0.8$ & $+8.6$ dB at \SI{1155}{Hz} & $+8.1$ & $+9.9$ dB at \SI{1345}{Hz} \\
$15^\circ$ & $+0.8$ & $+8.4$ dB at \SI{972}{Hz} & $+7.7$ & $+9.4$ dB \\
$30^\circ$ & $+0.7$ & $+7.2$ dB at \SI{917}{Hz} & $+5.8$ & $+8.3$ dB \\
$45^\circ$ & $+0.6$ & $+6.5$ dB at \SI{891}{Hz} & $+5.0$ & $+6.7$ dB at \SI{1068}{Hz} \\
$60^\circ$ & $+0.4$ & $+5.4$ dB at \SI{917}{Hz} & $+3.8$ & $+5.0$ dB \\
$75^\circ$ & $+0.2$ & $+3.6$ dB at \SI{1995}{Hz} & $+1.6$ & $+3.0$ dB \\
$90^\circ$ & $0.0$ & $+2.1$ dB at \SI{891}{Hz} & $+0.3$ & $+1.7$ dB at \SI{673}{Hz} \\\bottomrule
\end{tabularx}\par\smallskip
Useful band: \SIrange{125}{2000}{Hz} ($ka$ from 0.3 to 4.6). Below \SI{200}{Hz}
every angle is within \SI{1}{dB} of \SI{0}{dB}: the body is acoustically invisible. From
\SI{200}{Hz} the response rises, head-on the fastest, and up to \SI{2}{kHz} the measured
curves follow the BEM for the same angles within about \SI{1}{dB} and in the right
order ($0^\circ$ highest, $90^\circ$ lowest, within \SI{+2}{dB}). The $0^\circ$ maximum is
\SI{1.3}{dB} lower and 15\,\% lower in frequency than the BEM centre-of-face value, all
oblique curves fall away above \SI{2}{kHz}, and there are notches of 10--\SI{19}{dB}
between 2.5 and \SI{6}{kHz}. These deviations are explained in the next section.

\subsection{Position of the UMIK tip}\label{sec:tip}
\fig{labB_tip_position.png}{$0^\circ$: measured (solid) against the BEM at the face centre, at a field point \SI{3}{cm} in front of the face, and averaged over the face.}
The BEM evaluated \SI{3}{cm} in front of the face reproduces the measured $0^\circ$ curve
within \SI{1.5}{dB} from \SI{200}{Hz} to \SI{3.5}{kHz}: the same maximum ($+8.7$~dB at
1.1--\SI{1.2}{kHz}), the same dip to $-1.5$~dB at \SI{2.5}{kHz}, a notch of $-21$~dB at
\SI{4.0}{kHz} against the measured $-18$~dB at \SI{3.9}{kHz}, and the return to $+9$~dB
at \SI{5}{kHz}. A rigid face reflects the wave onto itself and a standing wave forms in
front of it; a quarter wavelength in front the incident and reflected waves cancel
($\lambda/4=\SI{2.2}{cm}$ at \SI{3.9}{kHz}), consistent with the recorded gap of
2--\SI{3}{cm}. At oblique incidence the path difference shrinks with $\cos\theta$ and the
notch moves up in frequency; the field-point BEM places the $15^\circ$, $30^\circ$ and
$45^\circ$ notches at 2.7, 3.4 and \SI{4.3}{kHz} (measured 2.5, 3.35 and \SI{3.9}{kHz}).
The notch depths at oblique angles do not agree, but the BEM uses 8 circumferential
terms and its oblique curves are not reliable above about \SI{4}{kHz}. The measured
data is therefore a centre-of-face free-field correction up to about \SI{2}{kHz} and a
measurement of the field \SIrange{2}{3}{cm} in front of the body above that.

\subsection{Reference stability (additional quality check)}
No end reference was taken, so $\Delta L_{\mathrm{drift}}$ cannot be evaluated. The
indirect evidence is that every normalized curve sits at $0.0$ to $+0.8$~dB between 125
and \SI{200}{Hz}, in the order and at the level the BEM predicts, which a moved
microphone or a changed gain would not give. Part~0 confirms that the gain did not
change between the reference and the distance runs.

\newpage
\section{Part 2: BEM simulation and spatial effects}
The supplied \file{CylinderPlaneWave} model assumes plane-wave incidence. Its frequency
grid is \SIrange{50}{5382}{\hertz} in $1/12$-octave steps. \file{AngleS} is in radians;
\file{CylEnd} selects a round or flat back end. Geometry used: $D=\SI{250}{mm}$,
$L=\SI{855}{mm}$, round back (\file{Lab B/bem/run_bem.m}, angles $0^\circ$ to
$180^\circ$); the field point is \SI{3}{cm} in front of the face centre on the axis.
Comparison is made only over the common band \SIrange{50}{5382}{Hz}.

\par\smallskip
\begin{tabularx}{\linewidth}{@{}lX@{}}\toprule
Output / choice & Result at $0^\circ$ \\\midrule
\file{p_centre} & $+9.9$~dB at \SI{1345}{Hz} ($ka=3.1$), dip to $-1.0$~dB at \SI{2.7}{kHz}, back to $+9.9$~dB at \SI{4}{kHz}. \\
\file{p_FP} (\SI{3}{cm} in front) & $+8.7$~dB at 1.1--\SI{1.2}{kHz}, $-21$~dB at \SI{4.0}{kHz}: the measured curve. \\
\file{p_avg} & $+7.8$~dB maximum, $+6.2$~dB at \SI{4}{kHz}: what a diaphragm of the full face diameter would feel; smoother at high $ka$. \\
\file{Round} vs. \file{Flat} & Less than \SI{0.25}{dB} difference at $0^\circ$: the back end only matters for sound from behind. \\\bottomrule
\end{tabularx}\par\smallskip
\fig{labB_bem_positions.png}{BEM, $0^\circ$: centre, field point \SI{3}{cm} in front, parabolic average over the face, and flat back end.}

\subsection{Agreement and limitations}
\par\smallskip
\begin{tabularx}{\linewidth}{@{}lXX@{}}\toprule
Comparison & Observation / quantitative difference & Possible explanation \\\midrule
Measured vs. face centre & Within \SI{1}{dB} up to \SI{2}{kHz} at every angle; $0^\circ$ maximum $-1.3$~dB and $-15$\,\% in frequency; notches above \SI{2.5}{kHz} & Tip not flush (2--\SI{3}{cm}); protractor $\pm2$--$3^\circ$; spherical incidence; nominal diameter \\
Measured vs. field point & Within \SI{1.5}{dB} up to \SI{3.5}{kHz} at $0^\circ$, notch frequency within 4\,\% & This is the pressure that was actually measured \\
Centre vs. average & Average \SI{2}{dB} lower at the maximum, \SI{3.7}{dB} lower at \SI{4}{kHz} & Pressure not uniform over the face at $ka>3$; a real diaphragm averages \\
Round vs. flat rear & $<\SI{0.25}{dB}$ at $0^\circ$ & Front face decides head-on \\\bottomrule
\end{tabularx}\par\smallskip
The field-point pressure best represents the experiment. Other contributions: the
UMIK body and its stand scatter too (the $\pm3$~dB ripples above \SI{400}{Hz} in Part~0
are already present without the mock-up and partly cancel in the ratio); the room below
\SI{125}{Hz}; the sound speed on the day (1\,\% per \SI{3}{\celsius} on the frequency
axis); and the numerical resolution of the BEM above \SI{4}{kHz} at oblique angles.

\newpage
\section{Part 2 continued: Scale to real microphones}
For a geometrically similar microphone of diameter $D_{\mathrm{mic}}$, define
\begin{equation}
s=\frac{D_{\mathrm{mock}}}{D_{\mathrm{mic}}},\qquad
f_{\mathrm{mic}}=s f_{\mathrm{mock}}.
\end{equation}
The smaller microphone shows corresponding features at \emph{higher} frequencies.
Conversely, a microphone datasheet frequency maps to the mock-up axis as
$f_{\mathrm{mock}}=f_{\mathrm{mic}}/s$. The normalized pressure level is unchanged
under ideal geometric scaling with the same medium.

\par\smallskip
\begin{tabularx}{\linewidth}{@{}lXXXX@{}}\toprule
Nominal size & Diameter used / mm & Scale factor $s$ & Measured $0^\circ$ maximum lands at & Handbook type \\\midrule
$1''$ & 23.77 & 10.5 & \SI{12.1}{kHz} & B\&K 4144/4145 \\
$1/2''$ & 12.70 & 19.7 & \SI{22.7}{kHz} & B\&K 4133/4134 \\
$1/4''$ & 6.35 & 39.4 & \SI{45}{kHz} & B\&K 4135/4136 \\
$1/8''$ & 3.175 & 78.7 & \SI{91}{kHz} & B\&K 4138 \\\bottomrule
\end{tabularx}\par\smallskip
$D_{\mathrm{mock}}=\SI{250}{mm}$ is the nominal value; the real microphone diameters
are the nominal body diameters (the $1''$ body is \SI{23.77}{mm}, not \SI{25.4}{mm}).
\fig{labB_scaled_to_microphones.png}{The measured curves on the frequency axis of each microphone size.}
For the $1''$ microphone the handbook free-field correction (B\&K 4145 in the brief)
peaks at about \SI{+10}{dB} between 12 and \SI{15}{kHz} at $0^\circ$ and falls towards
\SI{0}{dB} at $90^\circ$; the scaled measurement puts $+8.6$~dB at \SI{12}{kHz}, the
right place and slightly low for the reason given in Sec.~\ref{sec:tip}. The scaled BEM
centre value would be $+9.9$~dB at \SI{14}{kHz}.
\todo{Identify the exact handbook figure and whether the protection grid is fitted; add the scaled BEM curve to the figure if the group wants a three-way comparison.}

\section{Uncertainty and conclusion}
\textbf{Sources, ranked by their effect on the result:}
\begin{enumerate}
\item UMIK tip \SIrange{2}{3}{cm} in front of the face: changes the result qualitatively above \SI{2}{kHz} and lowers the $0^\circ$ maximum by \SI{1.3}{dB}. Systematic, explained by the BEM field point.
\item No start reference: drift cannot be quantified; the low-frequency levels ($0.0$ to $+0.8$~dB) bound it to below \SI{1}{dB}.
\item Mock-up diameter not measured: a few per cent on every frequency in the scaled comparison, nothing on the dB axis. For $s=D_{\mathrm{mock}}/D_{\mathrm{mic}}$,
\begin{equation}
\frac{u(s)}{s}=\sqrt{\left(\frac{u(D_{\mathrm{mock}})}{D_{\mathrm{mock}}}\right)^2+
\left(\frac{u(D_{\mathrm{mic}})}{D_{\mathrm{mic}}}\right)^2}.
\end{equation}
\item Angle read from a paper protractor ($\pm2$--$3^\circ$): the curves are 0.5--\SI{1}{dB} apart per $15^\circ$ near the maximum.
\item Spherical incidence from \SI{2.8}{m}, temperature not recorded, room below \SI{125}{Hz}: each below \SI{1}{dB} or outside the useful band.
\end{enumerate}
All normalized curves share one reference, so their errors are correlated: a reference
error moves every angle together and does not change the angle dependence.
\todo{Give numbers with a stated coverage if the group wants a numerical budget (rectangular distribution for the protractor and the diameter).}

\textbf{Conclusion.} Between \SI{125}{Hz} and \SI{2}{kHz} ($ka=0.3$ to $4.6$) the
measured pressure increase at the face of the mock-up is \SI{0}{dB} below $ka\approx0.5$,
rises head-on to $+8.6$~dB at \SI{1155}{Hz} and stays within $+2$~dB at grazing
incidence, in the order and within about \SI{1}{dB} of the level the BEM predicts for
the face centre. Above \SI{2}{kHz} the measurement is that of a point 2--\SI{3}{cm} in
front of the face, which the BEM field point reproduces within \SI{1.5}{dB} up to
\SI{3.5}{kHz} including the \SI{3.9}{kHz} notch. Scaled to a $1''$ microphone the
maximum lands at \SI{12}{kHz}, where the handbook correction has its $+10$~dB. Part~0
confirms free-field conditions to \SI{0.1}{dB} between 125 and \SI{400}{Hz}.

\section*{Sources and processing record}
\small
DTU 34870, \emph{Lab B -- Microphone scattering -- E26},
\file{34870_Lab_B_CylinderScattering_E2026.pdf}: pp. 1--5 (procedure, equations and
BEM), followed by the microphone handbook excerpts. Source: course vault / DTU Learn.

Repository workflow: raw files from the lab PC in \file{Lab B/data/Lab B master group 10/}
(unchanged); \file{Lab B/matlab/import_group_files.m} (frequency axis and metadata,
writes \file{labB_*.mat}); \file{Lab B/bem/run_bem.m}; \file{Lab B/matlab/process_labB.m}
with \file{D_mockup = 0.250} (nominal). Course measurement files: the 21 September 2026
re-upload (\texttt{34870 - Lab B (2).zip}).

\textbf{Analysis version / selected repeats / deviations:} processed 22 September 2026;
no repeats; \file{test.mat} excluded; reference \file{no_mockup_1.8m.mat} identified as
the \SI{2.8}{m} run by the $1/r$ check; no start reference, hence no drift figure.
\end{document}
